\documentclass[11pt]{article}
\usepackage[margin=1in]{geometry}
\usepackage{amsmath,amssymb,amsthm,mathtools}
\usepackage{graphicx}
\usepackage{microtype}
\usepackage[colorlinks=true,linkcolor=blue,citecolor=blue,urlcolor=blue]{hyperref}

\newtheorem{theorem}{Theorem}
\newtheorem{lemma}{Lemma}
\newcommand{\R}{\mathbb R}
\newcommand{\Z}{\mathbb Z}
\newcommand{\cM}{\mathcal M}
\newcommand{\Dmag}{\nabla+iA}

\title{Unbounded Magnetic Ground-State Multiplicity\\
by Products of Half-Flux Annuli}
\author{Candidate full-resolution packet for arXiv:math/0205088}
\date{12 August 2026}

\begin{document}
\maketitle

\noindent\textbf{Status.}
Candidate full resolution of the final open problem in
Lieb--Seiringer--Yngvason \cite{LSY}, subject to human validity and novelty
review.

\section{Source question}

The source is E.\ H.\ Lieb, R.\ Seiringer, and J.\ Yngvason,
\emph{Poincar\'e inequalities in punctured domains},
arXiv:math/0205088, later published in \emph{Annals of Mathematics}
158 (2003), 1067--1080 \cite{LSY}.  On printed page 1079, its final open
problem asks:

\begin{quote}
The dimension of the ground state manifold \(\cM_A^{p,q}\) can be bigger
than one. In the appendix we give an example for the case \(p=q=2\) where
its dimension is two. It would be interesting to construct examples (or
prove their existence) where the dimension can be arbitrarily large.
\end{quote}

\begin{center}
\includegraphics[width=0.96\linewidth]{figures/open_problem_crop.png}
\end{center}

For \(p=q=2\), the source identifies \(\cM_A^{2,2}\) with the ground
eigenspace of the Neumann realization of
\[
        H_A=-(\nabla+iA)^2
\]
on a bounded connected domain with the cone property.  Equivalently, it is
the minimizer space for
\[
 E_A^{2,2}=\inf_{0\ne u\in H^1(\Omega)}
 \frac{\|(\nabla+iA)u\|_{L^2(\Omega)}}{\|u\|_{L^2(\Omega)}}.
\]

\section{Proof intuition}

The source already supplies the elementary multiplicity-two building block:
an annulus carrying half a quantum of Aharonov--Bohm flux.  Angular Fourier
modes then see the shifted integers \(\ell-\tfrac12\), so precisely the two
modes \(\ell=0,1\) minimize the radial energy.  The upgrade is to take
Cartesian products.  Magnetic Neumann operators on a product add, their
ground energies add, and their ground eigenspaces tensor.  Thus \(k\)
half-flux annuli have ground multiplicity \(2^k\).

\section{Formal result}

\begin{theorem}\label{thm:main}
For every integer \(k\geq1\), there are a bounded connected open set
\(\Omega_k\subset\R^{2k}\) with the cone property and a vector field
\(A_k\in C^\infty(\overline{\Omega_k};\R^{2k})\) such that
\[
       \dim \cM_{A_k}^{2,2}=2^k.
\]
Consequently, the dimensions of the ground-state manifolds occurring in the
setup of \cite{LSY} are arbitrarily large.
\end{theorem}

\begin{proof}
Fix \(0<a<b<\infty\), and let
\[
 D=\{x=(x_1,x_2)\in\R^2:a<|x|<b\},\qquad
 A_0(x)=\frac{(-x_2,x_1)}{2|x|^2}.
\]
Thus, in polar coordinates, \(A_0=(2r)^{-1}e_\theta\).  It is smooth
and bounded on \(\overline D\).  Let \(H_0\) be the self-adjoint magnetic
Neumann operator associated with the closed quadratic form
\[
 q_0[u]=\int_D|(\nabla+iA_0)u|^2\,dx,\qquad
 \operatorname{dom}q_0=H^1(D).
\]

Expand \(u\) in angular modes as
\[
             u(r,\theta)=\sum_{\ell\in\Z}e^{-i\ell\theta}g_\ell(r).
\]
Orthogonality gives
\[
 q_0[u]
 =2\pi\sum_{\ell\in\Z}\int_a^b
 \left(|g_\ell'(r)|^2+
 \frac{(\ell-\tfrac12)^2}{r^2}|g_\ell(r)|^2\right)r\,dr.
 \tag{1}\label{eq:decomp}
\]
For \(\nu\geq0\), write
\[
 \mu(\nu)=
 \inf_{0\ne g\in H^1(a,b)}
 \frac{\displaystyle\int_a^b
   \left(|g'|^2+\frac{\nu^2}{r^2}|g|^2\right)r\,dr}
 {\displaystyle\int_a^b |g|^2r\,dr}.
\]
If \(0\leq\nu_1<\nu_2\), then for every nonzero \(g\),
\[
 \frac{Q_{\nu_2}[g]}{\|g\|^2_{L^2(r\,dr)}}
 \geq
 \frac{Q_{\nu_1}[g]}{\|g\|^2_{L^2(r\,dr)}}
 \frac{\nu_2^2-\nu_1^2}{b^2}.
\]
Taking infima shows that \(\mu(\nu)\) is strictly increasing:
\[
 \mu(\nu_2)\geq\mu(\nu_1)+
 \frac{\nu_2^2-\nu_1^2}{b^2}>\mu(\nu_1).
 \tag{2}\label{eq:strict}
\]
The smallest possible value of \(|\ell-\tfrac12|\) is \(\tfrac12\),
attained exactly for \(\ell=0,1\).  Equations
\eqref{eq:decomp}--\eqref{eq:strict} therefore show that the bottom of
\(\sigma(H_0)\) occurs in exactly these two angular sectors.  In each
sector, the lowest radial eigenvalue is simple by the standard
one-dimensional regular Sturm--Liouville theorem.  Hence
\[
                     \dim\ker(H_0-\mu(\tfrac12))=2.
 \tag{3}\label{eq:two}
\]

Now set
\[
 \Omega_k=D^k\subset(\R^2)^k=\R^{2k}
\]
and, writing \(x=(x^{(1)},\ldots,x^{(k)})\), define
\[
 A_k(x)=\bigl(A_0(x^{(1)}),\ldots,A_0(x^{(k)})\bigr).
\]
The domain \(\Omega_k\) is bounded, connected, and has the cone property
(indeed it is a bounded Lipschitz product domain), while \(A_k\) is smooth
and bounded on its closure.

The magnetic form on \(\Omega_k\) is the sum of the forms acting in the
individual variables.  Thus its Neumann operator is the tensor sum
\[
 H_k=\sum_{j=1}^k
 I\otimes\cdots\otimes I\otimes H_0\otimes
 I\otimes\cdots\otimes I.
\]
Because \(H_0\) has compact resolvent, an orthonormal eigenbasis for \(H_0\)
produces, by tensor products, an orthonormal eigenbasis for \(H_k\).
The corresponding eigenvalues are all sums
\(\lambda_{n_1}+\cdots+\lambda_{n_k}\).  Their least value is
\(k\mu(\tfrac12)\), and equality holds exactly when every tensor factor
lies in the two-dimensional space in \eqref{eq:two}.  Therefore
\[
 \ker\!\left(H_k-k\mu(\tfrac12)\right)
   =\ker(H_0-\mu(\tfrac12))^{\otimes k}
\]
and
\[
             \dim\cM_{A_k}^{2,2}
             =\dim\ker\!\left(H_k-k\mu(\tfrac12)\right)
             =2^k.
\]
Since \(2^k\to\infty\), the requested dimensions are arbitrarily large.
\end{proof}

\section{Verification and scope}

\begin{itemize}
\item \textbf{Hypotheses.}  The annulus stays a positive distance from the
origin, so the half-flux potential is smooth and bounded.  Finite products
are bounded, connected, and satisfy the cone property.
\item \textbf{Exactness of the building block.}  The quantitative inequality
\eqref{eq:strict} excludes every angular mode except \(\ell=0,1\);
Sturm--Liouville simplicity excludes hidden radial multiplicity.
\item \textbf{Exactness of the product count.}  The tensor-sum spectrum shows
that no excited factor can contribute at the bottom.  The multiplicity is
therefore exactly \(2^k\), not merely at least \(2^k\).
\item \textbf{Scope.}  The construction lets the ambient dimension grow as
\(2k\).  This answers the literal source request, which places no
fixed-dimension restriction.  It does not claim the stronger assertion that
ground multiplicity is unbounded while a particular \(n\), such as \(n=2\),
is held fixed.
\end{itemize}

\section{Bounded novelty check}

The run's registry, solution, attempt, and proof-gap indexes were searched by
arXiv id, title, and the phrases "ground state manifold" and "magnetic
ground-state multiplicity"; no duplicate answer was found.  A local
full-source search found later citations including the fractional magnetic
extension arXiv:2309.06919, but no discussion resolving this multiplicity
question.  Web searches used the exact paper title together with
"arbitrarily large", "ground state manifold", "multiplicity",
"tensor product", and "product of annuli", and also searched arXiv for
half-flux annular magnetic Neumann ground states.  No prior statement of this
product construction was found in that bounded search.  This is not a
definitive specialist-literature certification.

\section{Human-review recommendation}

Verify (i) that the source question is intended literally rather than with an
unstated fixed-\(n\) requirement, (ii) the product-domain cone-property claim,
and (iii) the tensor-sum identification for magnetic Neumann forms.  A
specialist citation search should precede any novelty claim beyond
"candidate full resolution".

\begin{thebibliography}{9}
\bibitem{LSY}
E.\ H.\ Lieb, R.\ Seiringer, and J.\ Yngvason,
\emph{Poincar\'e inequalities in punctured domains},
Ann. of Math. (2) \textbf{158} (2003), 1067--1080;
arXiv:math/0205088.
\end{thebibliography}

\end{document}
